\section{USAMO 1972}
        \begin{enumerate}
        	\item (\textit{USAMO 1972}) The symbols $(a,b,\ldots,g)$ and $[a,b,\ldots, g]$ denote the greatest common divisor and least common multiple, respectively, of the positive integers $a,b,\ldots, g$. For example, $(3,6,18)=3$ and $[6,15]=30$. Prove that


 $\frac{[a,b,c]^2}{[a,b][b,c][c,a]}=\frac{(a,b,c)^2}{(a,b)(b,c)(c,a)}$.

 

	\item (\textit{USAMO 1972}) A given tetrahedron $ABCD$ is isosceles, that is, $AB=CD, AC=BD, AD=BC$. Show that the faces of the tetrahedron are acute-angled triangles.


 

	\item (\textit{USAMO 1972}) A random number selector can only select one of the nine integers 1, 2, ..., 9, and it makes these selections with equal probability. Determine the probability that after $n$ selections ($n>1$), the product of the $n$ numbers selected will be divisible by 10.


 

	\item (\textit{USAMO 1972}) Let $R$ denote a non-negative rational number. Determine a fixed set of integers $a,b,c,d,e,f$, such that for \textit{every} choice of $R$, 


 $\left|\frac{aR^2+bR+c}{dR^2+eR+f}-\sqrt[3]{2}\right|<|R-\sqrt[3]{2}|$

 

	\item (\textit{USAMO 1972}) A given convex pentagon $ABCDE$ has the property that the area of each of the five triangles $ABC$, $BCD$, $CDE$, $DEA$, and $EAB$ is unity. Show that all pentagons with the above property have the same area, and calculate that area. Show, furthermore, that there are infinitely many non-congruent pentagons having the above area property.


 
                    \begin{figure}[H]
                        \centering
                        \includegraphics[width=0.5\textwidth]{Problem/USAMO_I_1972_5_Problem_Figure_1.png}
                    \end{figure}
                    


 

\end{enumerate}